\documentclass[11pt]{article}

\usepackage[a4paper,margin=29mm]{geometry}
\usepackage{amsmath,amssymb,amsthm,mathtools}
\usepackage[hidelinks]{hyperref}
\usepackage{microtype}

\newtheorem{theorem}{Theorem}[section]
\newtheorem{corollary}[theorem]{Corollary}
\newtheorem{proposition}[theorem]{Proposition}
\newtheorem{remark}[theorem]{Remark}
\newtheorem{problem}[theorem]{Open problem}

\newcommand{\R}{\mathbb R}
\newcommand{\N}{\mathbb N}
\newcommand{\Ioc}[2]{(#1,#2]}
\newcommand{\vol}{\operatorname{vol}}
\newcommand{\K}{\mathcal K}
\newcommand{\A}{\mathcal A}
\newcommand{\F}{\mathcal F}

\title{\textsc{Retracted}: Displacement--Overlap Conservation in Bicausal Transport\\
and a Critical Surplus Law for Suzuki's Zeta Channel}
\author{HGT Research Project}
\date{22 July 2026}

\begin{document}
\maketitle

\begin{abstract}
\textbf{Retracted on 22 July 2026.}
The interval-overlap identity and tilted-tent integral computed below are mathematically correct,
but the manuscript's central interpretation is false.  After duplicate seam capacity is charged
back, mixed future/past slack is exactly the original capacity-one quadrature deficit; total
apex-relative flank deficit is not free transport capacity.  A stronger formal countermodel now
shows that even an atomless source satisfying the sharp width-two cumulative band, with the full
decreasing flank left untouched, can violate Suzuki's completed-lag inequality because its
rising-flank excess grows with the apex while the reflected-polar/Gamma reserve is uniformly
bounded.  This document is retained only as a record of the failed route and must not be cited as
an RH advance.  Corrected results are recorded in
\texttt{SuzukiBicausalNoFreeSplice.lean} and
\texttt{SuzukiCriticalApexVacuumCapacityNoGo.lean}.
\end{abstract}

\section{Introduction}

Suzuki's screw-function framework replaces distributional Weil-form questions by concrete
continuous and operator-theoretic objects and proposes a limiting spectral construction for the
imaginary parts of the nontrivial zeta zeros \cite{Suzuki2026}.  In the rectangular sublane, the
prime wall is a positive atomic sample of a completed archimedean lag kernel.  This makes a
transport proof of Weil positivity conceivable: assign each Mangoldt atom to Lebesgue target mass
on a monotone branch of the positive kernel.

That program has a specific obstruction.  Increasing branches require future cells; decreasing
branches require past cells.  Each family separately forms a disjoint cumulative-mass partition,
but a moving maximum forces the two causal orientations to meet.  The unresolved question is not
whether the one-sided inequalities are valid.  It is whether they can be spliced without using
the same target capacity twice.

General transport-majorization theorems compare densities through controlled transport maps
\cite{MelbourneRoberto2021}.  They do not identify the surcharge created when the orientation of
one cumulative partition reverses.  Finite Guinand--Weil constructions give exact finite
dictionaries and archimedean tail control \cite{Groskin2026}, but do not provide this prime-side
splice.  The log-Haar capacity inequality
\[
 \sum_{n\le X}\frac{\Lambda(n)}n\le\log X
\]
is available unconditionally in the literature \cite{RamareZuniga2026}; it locates sharp
one-sided cells, but it does not by itself resolve their mixed overlap.

This paper contributes two exact laws.

First, the \emph{displacement--overlap conservation law} states that a forward displacement $h$
reappears with coefficient one as seam overlap.  This is independent of arithmetic and of the
shape of the sampled profile.  It converts vague ``double counting near the apex'' into an exact
necessary-and-sufficient integral surcharge.

Second, the \emph{bilateral tilted-tent surplus law} computes all continuum flank slack at once.
Its critical Suzuki specialization gives a strict comparison in the favorable direction:
continuum surplus exceeds the worst width-two seam charge.  This reverses the conclusion suggested
by small one-flank block constants.  The positive-transport route is not geometrically
under-capacitated; its missing theorem is source-specific arithmetic matching.

The paper does not assert RH.  It also does not call the entire Suzuki program a dead end.  It
settles one major yes/no question inside that program: the continuum geometry survives the exact
seam test, whereas a proof using only a two-unit cumulative capacity band cannot close it.

\section{Contiguous bicausal transport}

Let a positive atomic source have cumulative masses $F_j$.  A one-sided monotonicity argument
replaces atoms by intervals of the same length.  Abstracting from the arithmetic labels, suppose
the aggregate future target through a cutoff is
\[
 \mathsf L=\Ioc{a}{s+h},
\]
while the aggregate past target after the cutoff is
\[
 \mathsf R=\Ioc{s}{b}.
\]
The number $h\ge0$ is the relative displacement between the two copies of the cumulative
partition.  We assume $a\le s$ and $s+h\le b$, so both sides extend beyond the seam.

\begin{theorem}[Displacement--overlap conservation]\label{thm:overlap}
Under the preceding inequalities,
\begin{align}
 \mathsf L\cap\mathsf R&=\Ioc{s}{s+h},\label{eq:intersection}\\
 \mathsf L\cup\mathsf R&=\Ioc{a}{b}.\label{eq:union}
\end{align}
In particular,
\[
 \vol(\mathsf L\cap\mathsf R)=h.
\]
For every real profile $\Phi$ interval-integrable on $[a,b]$,
\begin{equation}\label{eq:integral-ledger}
 \int_a^{s+h}\Phi(t)\,dt+\int_s^b\Phi(t)\,dt
 =\int_a^b\Phi(t)\,dt+\int_s^{s+h}\Phi(t)\,dt.
\end{equation}
\end{theorem}

\begin{proof}
Membership in the intersection means
\[
 a<t\le s+h\quad\text{and}\quad s<t\le b.
\]
Since $a\le s$ and $s+h\le b$, this is equivalent to $s<t\le s+h$, proving
\eqref{eq:intersection}.  If $a<t\le b$, then either $t\le s+h$, placing $t$ in
$\mathsf L$, or $s+h<t$, placing it in $\mathsf R$.  This proves \eqref{eq:union}; the reverse
inclusions are immediate.  The volume statement follows from the length of a half-open interval.

For \eqref{eq:integral-ledger}, split the first integral at $s$ and the second at $s+h$:
\begin{align*}
 \int_a^{s+h}\Phi+\int_s^b\Phi
 &=\left(\int_a^s\Phi+\int_s^{s+h}\Phi\right)
   +\left(\int_s^{s+h}\Phi+\int_{s+h}^b\Phi\right)\\
 &=\int_a^b\Phi+\int_s^{s+h}\Phi.
\end{align*}
\end{proof}

\begin{corollary}[Exact splice criterion]\label{cor:criterion}
Suppose one-sided monotonicity bounds an atomic source by the two aggregate integrals on the
left-hand side of \eqref{eq:integral-ledger}.  That argument is capacity one precisely when its
accumulated one-sided slack dominates
\[
 \int_s^{s+h}\Phi(t)\,dt.
\]
\end{corollary}

The theorem is a conservation law rather than an upper estimate.  Enlarging one-sided cells does
not remove the reversal cost: the enlargement changes $h$, and the duplicate interval changes by
the same amount.

\section{Sharp and unconditional Mangoldt seams}

Set
\[
 F_n=2\sum_{d\le n}\frac{\Lambda(d)}d,
 \qquad \tau_n=2\log n.
\]
The sharp past and future log-Haar cells are
\[
 P_n=\Ioc{F_{n-1}}{F_n},
 \qquad Q_n=\Ioc{F_{n-1}+2}{F_n+2}.
\]
Below a cutoff $j$, the future cells form the translated cumulative interval ending at $F_j+2$;
above it, the past cells start at $F_j$.  Theorem~\ref{thm:overlap} therefore yields the exact seam
\begin{equation}\label{eq:sharp-seam}
 \Ioc{F_j}{F_j+2},
\end{equation}
of volume two.  No zero statement enters this geometry.

The repository also contains a fully unconditional wider transport derived from an elementary
Mertens bound.  If its radius is $R\ge0$, past coordinates are shifted by $-R$ and future
coordinates by $R+2$.  Their relative displacement is
\begin{equation}\label{eq:wide-width}
 W=(R+2)-(-R)=2+2R.
\end{equation}
The aggregate overlap is consequently
\[
 \Ioc{F_j-R}{F_j+R+2},
\]
again of exact width $W$.  Thus the unconditional widening is valid one-sided transport, but it
cannot be mistaken for a splice: it enlarges the seam surcharge exactly as much as it enlarges the
causal reach.

\section{Bilateral tilted-tent surplus}

For $\beta,L>0$ define
\[
 r_{\beta,L}(x)=e^{\beta x}\left(1-\frac{|x|}{L}\right),
 \qquad -L\le x\le L.
\]
The factor $e^{\beta x}$ is the half-density tilt and the triangular term is a rectangular
autocorrelation tent.  Suzuki lag is $t=2x$, so the normalized two-flank deficit is
\begin{equation}\label{eq:Ddef}
 D(\beta,L)=2\int_{-L}^L\bigl(1-r_{\beta,L}(x)\bigr)\,dx.
\end{equation}

\begin{theorem}[Bilateral tilted-tent transform]\label{thm:transform}
For all $\beta,L>0$,
\begin{align}
 \int_{-L}^L r_{\beta,L}(x)\,dx
 &=\frac{2(\cosh(\beta L)-1)}{L\beta^2},\label{eq:transform}\\
 D(\beta,L)
 &=4L-\frac{4(\cosh(\beta L)-1)}{L\beta^2}.\label{eq:Dclosed}
\end{align}
\end{theorem}

\begin{proof}
On the negative and positive halves use, respectively, the primitives
\begin{align*}
 A_-(x)&=e^{\beta x}\left(\frac{1+x/L}{\beta}
                  -\frac{1}{L\beta^2}\right),\\
 A_+(x)&=e^{\beta x}\left(\frac{1-x/L}{\beta}
                  +\frac{1}{L\beta^2}\right).
\end{align*}
Direct differentiation gives
\[
 A_-'(x)=e^{\beta x}(1+x/L),\qquad
 A_+'(x)=e^{\beta x}(1-x/L).
\]
Endpoint evaluation gives
\begin{align*}
 \int_{-L}^0e^{\beta x}(1+x/L)\,dx
 &=\frac{\beta L-1+e^{-\beta L}}{L\beta^2},\\
 \int_0^Le^{\beta x}(1-x/L)\,dx
 &=\frac{e^{\beta L}-\beta L-1}{L\beta^2}.
\end{align*}
After addition, the linear terms cancel and
$e^z+e^{-z}=2\cosh z$ proves \eqref{eq:transform}.  Since
$\int_{-L}^L1\,dx=2L$, equation \eqref{eq:Dclosed} follows from \eqref{eq:Ddef}.
\end{proof}

\begin{theorem}[Strict critical seam excess]\label{thm:critical}
At $(\beta,L)=(1/2,2)$,
\begin{equation}\label{eq:critical}
 D(1/2,2)=16-8\cosh(1)>2.
\end{equation}
\end{theorem}

\begin{proof}
The identity is Theorem~\ref{thm:transform}.  The explicit bounds
\[
 e<2.7182818286,\qquad e^{-1}<0.3678794412
\]
give $2\cosh(1)=e+e^{-1}<7/2$, hence $\cosh(1)<7/4$.  Therefore
$16-8\cosh(1)>2$.
\end{proof}

\section{Exact Suzuki specialization}

Let $\A_{L,a}(t)$ be Suzuki's positive rectangular sample channel.  At width two define
\[
 t_a=4a-4,\qquad H_a=2e^{a-1}.
\]
Here $t_a$ is the moving-tent apex lag and $H_a$ its height once the fixed small-lag tent is
disjoint from the moving support.

\begin{theorem}[Critical Suzuki profile identity]\label{thm:suzuki-profile}
For $a\ge3$ and $x\in[-2,2]$,
\begin{equation}\label{eq:suzuki-profile}
 \A_{2,a}(t_a+2x)=H_a r_{1/2,2}(x).
\end{equation}
\end{theorem}

\begin{proof}
For the stated range, $t_a+2x\ge0$ and its half-lag is at least two, so the fixed tent vanishes.
For the moving tent,
\[
 \frac{t_a+2x}{2}-(2a-2)=x,
\]
and its height is $2-|x|$.  The exponential factor splits as
\[
 e^{(t_a+2x)/4}=e^{a-1}e^{x/2}.
\]
Multiplication gives
\[
 e^{a-1}e^{x/2}(2-|x|)
 =2e^{a-1}e^{x/2}(1-|x|/2)=H_a r_{1/2,2}(x).
\]
\end{proof}

Define the actual continuum two-flank surplus
\[
 \mathsf S_a=2\int_{-2}^2\bigl(H_a-\A_{2,a}(t_a+2x)\bigr)\,dx.
\]

\begin{corollary}[Actual Suzuki critical surplus]\label{cor:suzuki-surplus}
For every $a\ge3$,
\begin{equation}\label{eq:suzuki-surplus}
 \mathsf S_a=H_a(16-8\cosh1)>2H_a.
\end{equation}
\end{corollary}

\begin{proof}
Insert \eqref{eq:suzuki-profile}, factor out $H_a>0$, and apply
Theorem~\ref{thm:critical}.
\end{proof}

The result is scale-exact.  Both the seam ceiling and the surplus grow linearly with the moving
apex height.  The strict margin is therefore not a finite-aperture artifact and is not obtained by
optimizing a vanishing local estimate.

\section{What the theorem decides---and what it does not}

There are now two unconditional conclusions.

\begin{enumerate}
\item A common-shift future/past transport necessarily double-uses target mass.  At sharp
Mangoldt displacement its exact surcharge is one width-two seam integral.
\item The full critical Suzuki tent has strictly more total continuum flank surplus than the
maximum possible value $2H_a$ of that seam integral.
\end{enumerate}

The first conclusion rules out any argument that silently unions the one-sided cell families.
The second rules out a geometric dead-end verdict: there is enough continuum capacity in the
correct channel.

Neither conclusion assigns actual Mangoldt atoms to the available flank surplus.  The cumulative
band
\[
 \tau_n-2\le F_{n-1}\le F_n\le\tau_n
\]
cannot force such an assignment by itself.  An extremal admissible partition may place a full
width-two atom at a strict apex and no source mass on its flanks.  The unique-apex obstruction then
prevents pointwise payment, and grouping creates no surplus because there are no flank atoms.  This
is a countermodel to capacity-only reasoning; it is not the actual von Mangoldt measure.

The terminal source-specific question is therefore exact and falsifiable.

\begin{problem}[Suzuki--Mangoldt critical seam matching]\label{prob:matching}
Construct disjoint target blocks for the actual Mangoldt atoms in the width-two critical support
such that the total one-sided monotonicity slack assigned to the mixed seam is at least its seam
integral.  Alternatively, classify a source-canonical transport class broad enough to contain all
such block assignments and prove that no member can satisfy this inequality.
\end{problem}

A solution of the first alternative plugs into the already established completed-lag domination
criterion; a class-complete negative solution establishes a genuine dead end for that precise
positive-transport realization.  Merely proving another local source surplus, or refuting one
cell choice, does neither.

\section{Lean formalization}

Every theorem stated in this paper has a dependency-closed Lean~4 proof.  The formal map is:

\begin{itemize}
\item \texttt{HTGeometry/BicausalDisplacementOverlap.lean}: exact intersection, union, seam volume,
and integral conservation;
\item \texttt{HTGeometry/TiltedTentTwoFlankSurplus.lean}: bilateral transform, closed surplus law,
critical identity, and strict excess;
\item \texttt{RHTower/Suzuki/SuzukiMangoldtBicausalSeamConservation.lean}: sharp width-two and fully
unconditional Mertens-width Mangoldt seam ledgers;
\item \texttt{RHTower/Suzuki/SuzukiCriticalTwoFlankGeometry.lean}: exact Suzuki profile specialization and
strict actual-channel continuum surplus.
\end{itemize}

The proofs use only Lean and mathlib foundations.  No declaration of RH, no project axiom, no
\texttt{sorry}, no numerical certificate, and no external prime estimate occurs in these modules.
The formalization separates universal geometry from the unresolved actual-prime matching in
Problem~\ref{prob:matching}; the latter is not smuggled in as a theorem hypothesis.

\section{Conclusion}

Displacement--overlap conservation identifies the exact price of reversing a cumulative positive
transport: the displacement is returned as duplicate target capacity.  In Suzuki's sharp
Mangoldt geometry that price has width two.  The bilateral transform gives a second universal law
and shows that the critical width-two Suzuki channel carries strictly more total flank surplus
than this worst seam charge.

The combined verdict is sharper than either optimism or a no-go slogan.  Capacity-only cell
transport is insufficient; Suzuki continuum geometry is sufficient.  Remaining frontier is one
source-specific matching theorem for the actual Mangoldt measure.  That theorem, rather than
further one-flank estimates, is the next legitimate point at which this lane can reach RH or a
class-complete dead end.

\begin{thebibliography}{9}

\bibitem{Suzuki2026}
M.~Suzuki,
\emph{Weil's quadratic form via the screw function},
arXiv:2606.09096, 2026.

\bibitem{Groskin2026}
A.~Groskin,
\emph{A finite Guinand--Weil dictionary and archimedean tail order for the truncated Weil
quadratic form},
arXiv:2607.02828, 2026.

\bibitem{MelbourneRoberto2021}
J.~Melbourne and C.~Roberto,
\emph{Transport-majorization to analytic and geometric inequalities},
arXiv:2110.03641, 2021; revised 2022.

\bibitem{RamareZuniga2026}
O.~Ramar\'e and S.~Zuniga Alterman,
\emph{M\"obius function and primes: an identity factory with applications},
arXiv:2312.05138v3, 2026.

\end{thebibliography}

\end{document}
